% !TEX root = ../main.tex
%% Pages « Maîtrise de la langue » (édition 2027-2028)
%% Source : note de service du 11 septembre 2026, NOR MENE2622642N (BO spécial n° 4 du 17 septembre 2026),
%% épreuve terminale de l'enseignement de spécialité mathématiques, et son annexe « Attendus et observables rédactionnels ».
%% Les libellés de la grille sont repris de l'annexe officielle. La conversion en points est une proposition OptiTAB.

\thispagestyle{otFront}
\phantomsection\label{maitrise-langue}%
\pdfbookmark[1]{Maîtrise de la langue}{bm-maitrise-langue}%
\otPageTitre{La maîtrise de la langue}{Deux points sur vingt à l'épreuve écrite de spécialité, à partir de la session 2027}

{\fontsize{10.5}{14.5}\selectfont
\otIntertitre{Ce que dit le texte officiel}
La note de service du 11 septembre 2026 qui définit l'épreuve précise que « la maîtrise de la langue est prise en compte à
hauteur de deux points sur vingt, dédiés à la maîtrise des normes orthographiques et syntaxiques, ainsi qu'à la capacité à
formuler un raisonnement et à utiliser un vocabulaire juste et adapté ». Les attendus sont détaillés dans une grille annexée,
reproduite page suivante. La structure de l'épreuve ne change pas : quatre heures, quatre exercices indépendants, chacun noté
de $4$ à $8$ points.\par\vspace{3mm}
Une note de service du même jour, applicable à tous les examens à partir de la session 2027, demande aussi de « faire
davantage rédiger les candidats » en limitant les questionnaires à choix multiples et les réponses courtes. Dans ce livre, les
questionnaires des sujets type bac demandent donc toujours une justification rédigée.\par}

\vspace{7mm}
\noindent
\begin{minipage}[t]{0.48\linewidth}\vspace{0pt}
\begin{tcolorbox}[enhanced,colback=otNavyTint,colframe=otNavyTint,boxrule=0pt,arc=2mm,left=4.5mm,right=4.5mm,top=3.5mm,
  bottom=3.5mm,fontupper=\fontsize{10}{13.5}\selectfont]
\otIntertitre{En pratique, sur ta copie}
\begin{itemize}[label={\color{otNavy}\textbullet},leftmargin=1.1em,itemsep=1.4mm,topsep=0mm]
  \item Fais des phrases : « Comme\dots, on a\dots, donc\dots ».
  \item Annonce la propriété utilisée avant de l'appliquer.
  \item Emploie le mot juste : \emph{croissante}, \emph{majorée}, \emph{converge}, \emph{orthogonal}\dots
  \item Termine chaque question par une conclusion qui répond à la question posée.
  \item Relis-toi : accords, conjugaison, ponctuation.
\end{itemize}
\end{tcolorbox}
\end{minipage}\hfill
\begin{minipage}[t]{0.48\linewidth}\vspace{0pt}
\begin{tcolorbox}[enhanced,colback=otGreenTint,colframe=otGreenTint,boxrule=0pt,arc=2mm,left=4.5mm,right=4.5mm,top=3.5mm,
  bottom=3.5mm,fontupper=\fontsize{10}{13.5}\selectfont]
\otIntertitre{S'auto-évaluer avec les corrigés}
Les corrigés des sujets détaillent le barème des exercices sur $18$ points. Pour les $2$ points de langue, place ta copie dans
la grille, ligne par ligne. Le texte officiel ne fixe pas de conversion en points ; à titre indicatif, OptiTAB propose de
compter $0{,}5$ point par ligne évaluée « satisfaisant » ou « très satisfaisant », et $0$ sinon.
\end{tcolorbox}
\end{minipage}

\vspace{7mm}
\otIntertitre{Deux rédactions d'une même réponse}
{\fontsize{10}{13.5}\selectfont
\emph{Question :} démontrer que la fonction $f$ définie sur $\R$ par $f(x)=x^{3}+x$ est strictement croissante sur $\R$.\par}
\vspace{3mm}
\noindent
\begin{minipage}[t]{0.42\linewidth}\vspace{0pt}
\begin{tcolorbox}[enhanced,colback=otOrange!10,colframe=otOrange!10,boxrule=0pt,arc=2mm,left=4mm,right=4mm,top=3mm,bottom=3mm,
  fontupper=\fontsize{10}{13.5}\selectfont]
\otEtiq{otOrange!75!black}{À éviter}\par\vspace{2mm}
{\itshape $f'=3x^{2}+1>0$ $\Rightarrow$ croissante}\par\vspace{2mm}
{\fontsize{8.6}{11.5}\selectfont\color{otMuted}Pas de phrase, symbole $\Rightarrow$ employé pour « donc », justification du signe
absente, conclusion incomplète.\par}
\end{tcolorbox}
\end{minipage}\hfill
\begin{minipage}[t]{0.54\linewidth}\vspace{0pt}
\begin{tcolorbox}[enhanced,colback=otGreenTint,colframe=otGreenTint,boxrule=0pt,arc=2mm,left=4mm,right=4mm,top=3mm,bottom=3mm,
  fontupper=\fontsize{10}{13.5}\selectfont]
\otEtiq{otGreenDark}{Rédaction attendue}\par\vspace{2mm}
La fonction $f$ est dérivable sur $\R$ et, pour tout réel $x$, $f'(x)=3x^{2}+1$. Comme $x^{2}\geq0$, on a $f'(x)\geq1>0$. La
dérivée de $f$ est strictement positive sur $\R$, donc $f$ est strictement croissante sur $\R$.
\end{tcolorbox}
\end{minipage}

\vspace{6mm}
{\fontsize{10}{13.5}\selectfont\color{otMuted}
Pour t'entraîner à rédiger, le chapitre 23 propose des exercices de logique et de rédaction : négations, implications,
raisonnements, et une copie à corriger (exercice 12).\par}
\clearpage

%% ---------------------------------------------------------------------
\thispagestyle{otFront}
\phantomsection\label{grille-langue}%
\otPageTitre{La grille officielle}{Attendus et observables rédactionnels, annexés à la note de service du 11 septembre 2026}

{\fontsize{9.4}{12.2}\selectfont\renewcommand{\arraystretch}{1.5}\arrayrulecolor{otRule}%
\begin{tabularx}{\linewidth}{@{}>{\leavevmode\color{otNavy}\otSemi\fontsize{8.6}{11}\selectfont\raggedright\arraybackslash}p{24mm}YYYY@{}}
\toprule
 & \otTete{Très insuffisant} & \otTete{Insuffisant} & \otTete{Satisfaisant} & \otTete{Très satisfaisant} \\
\midrule
Orthographe & Erreurs très fréquentes, notamment sur des règles simples (accords sujet-verbe, accords dans le groupe nominal, etc.).
  & Erreurs fréquentes et/ou répétées (accords, conjugaison, homophones, etc.).
  & Erreurs ponctuelles, lexicales ou grammaticales (accords, homophones grammaticaux, etc.), mais qui ne se répètent pas de façon
    systématique.
  & Orthographe lexicale et grammaticale maîtrisée ; erreurs rares, isolées, sans caractère systématique. \\
\midrule
Syntaxe, construction des phrases & Syntaxe très défaillante : nombreuses phrases mal construites, incorrectes, incompréhensibles
    ou à la ponctuation inopérante.
  & Construction de phrases fragile : maladresses récurrentes, phrases parfois mal structurées.
  & Phrases globalement bien construites malgré quelques maladresses syntaxiques ponctuelles.
  & Phrases correctement construites, syntaxe maîtrisée voire riche (subordination, coordination, etc.). \\
\midrule
Lexique & Le lexique, notamment celui de la discipline, est employé à mauvais escient ou très pauvre.
  & Le lexique, notamment celui de la discipline, est pauvre et parfois employé à mauvais escient.
  & Le lexique, notamment celui de la discipline, est globalement correct.
  & Le lexique, notamment celui de la discipline, est riche et correctement utilisé. \\
\midrule
Mise en forme et organisation de la réflexion & Le propos n'est pas organisé.
  & Le propos est organisé de manière peu pertinente.
  & Le propos est organisé de manière globalement cohérente.
  & Le propos est organisé de manière cohérente. Il suit un fil conducteur perceptible et pertinent. \\
\bottomrule
\end{tabularx}}
\par\vspace{3mm}
{\fontsize{8}{10.5}\selectfont\color{otMuted}Source : annexe « Attendus et observables rédactionnels » de la note de service du
11 septembre 2026, NOR MENE2622642N, Bulletin officiel spécial n\textsuperscript{o}\,4 du 17 septembre 2026. Les mêmes paliers
figurent dans la note de service NOR MENE2623195N, applicable à tous les examens.\par}

\clearpage
